% Teaching piecewise-linear curve, not measured data.
% (episode,return): (0,.9),(20,.9),(20,.3),(25,.3),(30,.55),
% (35,.75),(40,.85),(50,.9),(60,.9).
% Change=20, detection=25, recovery=40 at predeclared threshold .85.
\begin{tikzpicture}
 \begin{axis}[width=.98\linewidth,height=67mm,
   font=\figurefont\fontsize{8.5}{11}\selectfont,
   axis lines=left,xmin=0,xmax=60,ymin=0,ymax=1.18,
   xtick={0,20,25,40,60},ytick={0,.3,.6,.9},
   xlabel={部署回合$k$},ylabel={假设回报},
   tick label style={font=\figurefont\fontsize{8}{10}\selectfont}]
  \path[fill=BookAmber!12] (axis cs:20,.9) -- (axis cs:40,.9)
    -- (axis cs:40,.85) -- (axis cs:35,.75) -- (axis cs:30,.55)
    -- (axis cs:25,.3) -- (axis cs:20,.3) -- cycle;
  \draw[BookRule,dashed] (axis cs:0,.85) -- (axis cs:60,.85);
  \draw[BookAmber,densely dashed] (axis cs:20,0) -- (axis cs:20,1.05);
  \draw[BookInk,dotted,line width=.8pt] (axis cs:25,0) -- (axis cs:25,.98);
  \draw[BookTeal,dash dot] (axis cs:40,0) -- (axis cs:40,1.05);
  \addplot[BookInk,line width=1pt] coordinates
    {(0,.9)(20,.9)(20,.3)(25,.3)(30,.55)(35,.75)(40,.85)(50,.9)(60,.9)};
  \node[anchor=east] at (axis cs:20,1.09) {变化};
  \node[anchor=west] at (axis cs:25,1.00) {检测};
  \node[anchor=south] at (axis cs:40,1.06) {恢复};
  \node[anchor=west,align=left] at (axis cs:31,.36)
    {检测延迟：5回合\\检测后恢复：15回合};
  \node[anchor=west] at (axis cs:1,.77) {恢复阈值$0.85$};
  \draw[{Latex[length=1.6mm]}-{Latex[length=1.6mm]},BookAmber]
    (axis cs:17,.3) -- (axis cs:17,.9);
  \node[anchor=east,align=right] at (axis cs:16,.48) {性能\\下降};
 \end{axis}
\end{tikzpicture}
